\section{A deepest weakly essential residual system}
\label{a:essential}

We return to the sequence of anticanonical cones \(W=\Spec\mathcal A\)
from Section~\ref{a:cone}, retaining its fixed integer \(n>1\), degree
action \(D\), canonical generator \(\Omega\), and vertex \(o\).
Constants called uniform in this section are independent of the member
of that sequence. The small parameters realizing each exact minimum may
depend on the individual member.

\subsection{Initial minimizing data}

Apply Proposition~\ref{a:exact-penalization} on \(W\) with constraint
\(I_n^{1/n}\) and residual discrepancy \(a=A_W\). Proposition
\ref{a:initial-lc-place} supplies a centred quasi-monomial lc place.
Normalize the resulting \(D\)-invariant minimizer \(\tilde v\) by
\(A_W(\tilde v)=1\). In constructing its auxiliary klt pairs, take
the general core equations in degree \(n\); every such equation has
value \(n\) at \(\tilde v\). On both \(\tilde v\) and the lc place
\(u^0\) from \eqref{eq:A2}, the auxiliary discrepancy is the parameter
times bare discrepancy, up to their common scaling. Here rescale
\(u^0\) so that \(A_W(u^0)=1\), which preserves \eqref{eq:A2}. Applying
\eqref{eq:A3} with minimizer \(\tilde v\) and test valuation \(u^0\)
therefore gives, uniformly on the ideal \(\m_z\),
\begin{equation}
                         \tilde v(\m_z)\longrightarrow\infty.
 \label{eq:A7}
\end{equation}
Indeed \(A_W(u^0)=A_W(\tilde v)=1\) and
\(u^0(\m_z)\to\infty\) by \eqref{eq:A2}.

Stable degeneration makes \(R=\gr_{\tilde v}\mathcal A\) finitely
generated and normal. Removing its effective specialized boundary shows
that \(\Spec R\) is klt. These homogeneous assertions, initially near
the vertex, hold globally: the positive grading makes every orbit
closure contain the vertex. Lemma~\ref{a:canonical-specialization}
gives the canonical generator \(\Omega_{\tilde v}\). Let \(T\) be the
joint effective torus of the value grading and \(D\), with character
lattice \(M_T\), and let \(\chi\) be the character of this form. Then
\[
                      D(\chi)=1=\tilde v(\chi).
\]
Every specialized core equation has pure weight \(n\chi\): its two
weights are degree \(n\) and value \(n\), and these determine the
joint character. Write \(I_R=(R_{n\chi})\).

The exact minimum fixes the normalized valuation and the chosen core
through all sufficiently small rational parameters on this member.
The induced auxiliary klt discrepancies, before their common canonical
scaling, are the specialized core discrepancy plus the parameter times
the bare discrepancy. Letting that parameter tend to zero shows that
the specialized core boundary is lc. The order of \(I_R\) is no larger
than the averaged order of the core equations; hence \(I_R^{1/n}\) is
lc as well. Moreover every quasi-monomial lc place \(u\) of this system
centred at the vertex satisfies
\begin{equation}
             u(f)\leq C A_R(u)\tilde v(f)
                        \qquad(0\ne f\in R),
 \label{eq:A8}
\end{equation}
with uniform \(C\). Apply \eqref{eq:A3} on the auxiliary klt stable
cone. At \(u\), the core equations have orders at least \(u(I_R)\),
so the auxiliary discrepancy is at most the parameter times \(A_R(u)\).
At the minimizing grading \(\tilde v\) it is exactly the parameter,
by \eqref{eq:A6}. The parameter cancels, giving the bound.

We may ensure that evaluation by \(\tilde v\) is injective on \(M_T\).
If necessary add a sufficiently small positive multiple \(\delta D\)
and renormalize by \(1+\delta\). On each original homogeneous degree
this translates values by \(\delta\) times that degree; extend by the
minimum rule. Generic \(\delta\) avoids ties of distinct joint weights
and retains the same joint graded ring and Abhyankar equality. To check
canonical normalization, choose homogeneous lifts of joint homogeneous
independent initial probes. In their logarithmic basis the coefficient
of \(\Omega\) has degree one and old value one, so its new value is
\(1+\delta\). Formula~\eqref{eq:A6} gives both the stated normalization
and the unchanged leading form. Homogeneous lifts exist by character
splitting in the \(D\)-invariant filtrations, also for fractions.

The actual ideal \(I_n\) still has normalized value \(n\).
All degrees on \(\mathcal A\) are nonnegative, and \(\m_z\) has degree
zero. Taking \(\delta\) individually small, the bounds
\eqref{eq:A7}--\eqref{eq:A8} persist with uniform factors. We henceforth
use the adjusted notation. To compare torus ranks, we will retain actual
valuations on \(W\), together with
\[
 \begin{gathered}
 A_W(\tilde v)=1,\quad \tilde v(I_n)=n,\quad
 R=\gr_{\tilde v}\mathcal A\ \text{finitely generated and klt},\\
 T\supset D\ \text{the value torus},\quad
 D(\chi)=\tilde v(\chi)=1,\quad I_R^{1/n}\ \text{lc},
 \end{gathered}
\]
and the uniform estimates \eqref{eq:A7}--\eqref{eq:A8}.

Among all sequences of such data and their subsequences, choose one
whose torus rank is constant and as large as possible. This is possible
because \(\rank T\leq\dim R=d+1\). All future rank comparisons concern
sequences retaining the uniform estimates, rather than isolated members.
For each member choose \(h\in\mathcal A_n\) general so that its divisor
divided by \(n\) is lc on \(W\) and the corresponding divisor is lc on
\(X\), and its initial \(\bar h\) is general nonzero in \(R_{n\chi}\).
Indeed this finite-dimensional weight space is exactly the leading
space at the minimal value \(n\) from degree \(n\). Include lifts of
its basis, together with system generators, in the finite-dimensional
space used for the general choice. On \(R\), the divisor
\(\Div(\bar h)/n\) is lc, and each of its discrepancy-zero primes is
also an lc place of \(I_R^{1/n}\). On a resolution, the moving part has
coefficient \(1/n<1\), so its general smooth components introduce no
additional discrepancy-zero primes. The fixed part is the ideal's part.

\subsection{The projective torus quotient}
\label{a:quotient}

The character \(\chi\) lies in the interior of the rational polyhedral
cone of occurring weights. Otherwise a nonzero rational dual
cocharacter would be nonnegative on all occurring weights and nonpositive
on \(\chi\). It defines a nontrivial centred valuation on \(R\), whose
bare discrepancy is its evaluation on \(\chi\), contradicting klt.
Also \(\chi\) is primitive, since \(D(\chi)=1\). Set
\[
 R_\chi=\bigoplus_{m\geq0}R_{m\chi},\qquad
 S=\Proj R_\chi.
\]
This ray ring is a subtorus invariant ring, hence finitely generated and
normal. Its field and the invariant field satisfy
\[
           k(S)=k(R)^T,\qquad \Frac R_\chi=k(S)(z^\chi).
\]
Here and below \(z^\mu\) are multiplicative rational torus variables.
We give the monoid detail that ensures these identities in all degrees.
The finitely generated occurring monoid spans \(M_T\) as a group.
Its saturation in the weight cone is a finite module over it: subtract
integer multiples of its generators until the remainder lies in a
bounded set. For each of the finitely many module representatives, use
the group-span property to write its difference as a difference of
monoid elements. Adding a common monoid translate sends the entire
saturation into the monoid. Since \(\chi\) is interior, for every
fixed weight \(\mu\) the points \(m\chi-\mu\) lie in that translated
saturation for all sufficiently large \(m\). Thus a homogeneous
numerator and denominator of equal weight can be multiplied by a common
occurring element to have ray degrees. Likewise ray elements occur in
all sufficiently large integral degrees, so their degree group is
\(\Z\chi\). This proves the field identities.

Choose a rational volume form \(\omega_S\) so that
\[
               \Omega_{\tilde v}
                    =z^\chi\omega_S\wedge d\log\mathbf z,
\]
where \(d\log\mathbf z\) is the wedge for a lattice basis of \(M_T\).
The quotient of the two sides is invariant, so it can be absorbed in
\(\omega_S\). There is an ample rational \(\Q\)-Cartier divisor \(L\)
on the normal projective variety \(S\) for which
\begin{equation}
 fz^{m\chi}\in R
       \quad\Longleftrightarrow\quad \Div_S f+mL\geq0
             \qquad(m\geq0,\ f\in k(S)^*).
 \label{eq:A9}
\end{equation}
Take a sufficiently divisible Veronese of the ray ring and write its
tautological polarization using the corresponding power of the rational
section \(z^\chi\). The equivalence holds in sufficiently high divisible
degrees by the Proj description. For any remaining degree, take a
sufficiently divisible power: membership of that power implies
membership of the element by normality of \(R_\chi\), and the divisor
inequality is equivalent to that for its power. If \(\dim S=0\), set
\(L=0\); all subsequent prime tests on \(S\) are then vacuous.

On \(S\) define
\[
 f_h=\bar h/z^{n\chi},\qquad
 H=L+\tfrac1n\Div_S f_h,\qquad
 B=-K_S^{\omega_S}-L.
\]
In particular \(H\geq0\) and \(H\sim_{\Q}L\).

\begin{lemma}[Gauss lifts and the quotient boundary]\label{a:gauss-lifts}
For every prime-divisor valuation \(P\) over \(S\), there is a rational
Gaussian extension to \(k(R)\), centred on \(\Spec R\), whose restriction
is \(P\) and whose value on \(z^\chi\) is \(P(L)\). It is divisorial
up to scaling. The pair \((S,B)\) is klt log Fano, and \(B\) is effective
with coefficients in \(\{1-1/j:j\in\Z_{>0}\}\).
\end{lemma}

\begin{proof}
Impose nonnegativity on finitely many eigenfunction generators of \(R\).
For an extension of \(P\) these are linear inequalities in the rational
torus-variable weights, with the one weight on \(z^\chi\) fixed to
\(P(L)\). The criterion for a solution is that every nonnegative
combination eliminating the other weights has nonnegative right-hand
side. By rational polyhedrality it suffices to test rational
combinations. Clear denominators and multiply the corresponding
eigenfunctions. Their total weight is a multiple of \(\chi\); it is a
nonnegative multiple because the occurring cone is pointed. The required
inequality is exactly \eqref{eq:A9} pulled back to \(P\). The rational
linear system therefore has a rational solution. Its Gaussian valuation
has rational rank one and satisfies Abhyankar equality, so it is
divisorial up to scaling. Formula~\eqref{eq:A6} gives its discrepancy
\[
                  A_B(P)=a_{\omega_S}(P)+P(L)>0,
\]
where positivity follows from klt of \(R\). Thus \((S,B)\) is sub-klt;
we next verify effectivity and the coefficients.

For a strict prime \(P\) on \(S\), first lift \(P\) to a homogeneous
prime divisor \(P'\) of the affine ray cone. Use the punctured cone for
a degree-one generated Veronese, and then the finite map from the full
ray cone. The affine subtorus quotient \(\Spec R\to\Spec R_\chi\) is
surjective: projection to invariant summands splits the inclusion as an
\(R_\chi\)-module. Localize at the generic point of \(P'\) and pull
back its uniformizer. A height-one minimal prime of this nonzero
parameter ideal gives a prime divisor \(E\) on \(\Spec R\) dominating
\(P'\). Its contraction is the maximal ideal of that discrete valuation
ring. A connected torus preserves each of these finitely many components,
so \(E\) is invariant. Its restriction is \(jP\) for an integer
\(j>0\), by Lemma~\ref{a:abhyankar}.

Choose a section in a sufficiently divisible degree whose corresponding
section of \(L\) is nonvanishing at the generic point of \(P\).
It has order zero at the ray-cone divisor and at its dominating lift,
which gives \(E(z^\chi)=jP(L)\). As \(E\) is a strict prime of the
variety \(\Spec R\) with zero canonical divisor, \(A_R(E)=1\).
The leading-form formula now yields
\[
                         1=jA_B(P).
\]
The coefficient of \(B\) at \(P\) is therefore \(1-1/j\), proving
effectivity and the asserted standard coefficients. Together with all
prime tests already proved, this makes \((S,B)\) klt. Finally
\(-K_S-B=L\) is ample, so it is log Fano.
\end{proof}

Test the lc divisor \(\Div(\bar h)/n\) on the Gauss lifts of
Lemma~\ref{a:gauss-lifts}. This gives, for every prime over \(S\),
\begin{equation}
 A^+(P):=A_B(P)-P(H)
       =a_{\omega_S}(P)-\tfrac1n P(f_h)\geq0.
 \label{eq:A10}
\end{equation}
The positive values of \(A^+\) on primes have gap \(1/n\): form orders
and function orders are integers. These are also generalized Calabi--Yau
data with effective boundary \(B+H\), principal nef b-divisor
\(-\frac1n\Div(f_h)\), and total adjoint divisor literally zero. The
principal divisor here is taken on every model. It is rational Cartier
and numerically trivial, so is b-nef. Indeed
\(K_S+B+H=\frac1n\Div_S f_h\).

For an invariant prime divisor \(E\) on \(\Spec R\), write
\(h_E=E(\bar h)/n\geq0\). When its restriction is nonzero,
Lemma~\ref{a:abhyankar} writes it as \(jP\), \(j\in\Z_{>0}\).
Then
\begin{equation}
 E|_{k(S)}=jP\ne0
       \quad\Longrightarrow\quad
       1=jA^+(P)+h_E,\qquad jP(H)\leq h_E.
 \label{eq:A11}
\end{equation}
For the equality, Gaussianity and \eqref{eq:A6} give
\[
 1=j a_{\omega_S}(P)+E(z^\chi),\qquad
 h_E=\tfrac jnP(f_h)+E(z^\chi).
\]
For the inequality, choose sufficiently divisible sections of \(L\)
avoiding the centre of \(P\). Their nonnegative values under \(E\),
together with \eqref{eq:A9}, give \(E(z^\chi)\geq jP(L)\).
If the restriction is trivial, the valuation is a pure toric
cocharacter and \(h_E=E(\chi)=A_R(E)=1\).

\subsection{Weak essentiality by maximality of the torus rank}

\begin{proposition}[Weak essentiality]\label{a:weak-essentiality}
On the chosen maximal-rank sequence, there are numbers \(\eta\to0\)
such that, for every prime divisor \(P\) over \(S\) with \(A^+(P)=0\),
\begin{equation}
       \sup_{H'\geq0,\ H'\sim_{\Q}H} P(H')
                            \leq(1+\eta)P(H).
 \label{eq:A12}
\end{equation}
\end{proposition}

\begin{proof}
Otherwise pass to a subsequence carrying lc primes \(P\) and effective
divisors \(H'\sim_{\Q}H\) with
\(P(H')>(1+c_*)P(H)\), for a fixed \(c_*>0\).
Here \(P(H)=A_B(P)>0\). The lift from
Lemma~\ref{a:gauss-lifts} has bare discrepancy \(P(H)\) and discrepancy
zero for \(\Div(\bar h)/n\). Normalize it to bare discrepancy one.
By our general choice of \(h\), it is also an lc place of \(I_R^{1/n}\).
Average this normalized lift and the grading valuation \(\tilde v\)
on eigenfunctions, extending by the Gaussian minimum rule. The resulting
valuation \(u_*\) restricts to \(P/(2P(H))\) on the invariant field.
It is quasi-monomial, has bare discrepancy one by \eqref{eq:A6}, and
is an lc place of \(I_R^{1/n}\). For the last assertion, the core
weight is pure: every core function has value at least \(n\) at the
lift and exactly \(n\) at the grading, with equality for \(\bar h\).
Also \(u_*\geq\tilde v/2\) on regular functions, so it is centred at
the vertex and
\[
                       \vol(u_*)\leq2^p\vol(\tilde v).
\]

Write \(H'=L+\frac1m\Div_S f'\) with \(m>0\) integral, clearing
rational denominators. By \eqref{eq:A9}, \(g=f'z^{m\chi}\in R_{m\chi}\).
Its value divided by \(m\) is \(P(H')/P(H)>1+c_*\) at the normalized
lift, and one at the grading. Hence
\[
                         u_*(g)/m>1+c_*/2.
\]

Choose rational \(1<c_0<1+c_*/2\) and a fixed rational \(\lambda>0\)
large enough that
\[
                         1+\lambda(c_0-1)>2C,
\]
where \(C\) is the old uniform constant in \eqref{eq:A8}.
On \(\Spec R\), minimize subject to the same lc constraint, now with
residual discrepancy
\[
 a(u)=(1+\lambda c_0)A_R(u)
          -\lambda\min\{u(g)/m,c_0u(I_R)/n\}.
\]
This has the ideal form of Proposition~\ref{a:exact-penalization}.
Explicitly take an integer \(q_0\) for which \(q_0/m\) and
\(q_0c_0/n\) are integral and put
\(J=(g^{q_0/m})+I_R^{q_0c_0/n}\); the subtracted term is
\((\lambda/q_0)u(J)\). Since the constraint ideal is lc,
\(u(I_R)/n\leq A_R(u)\), and consequently
\[
                  A_R(u)\leq a(u)\leq(1+\lambda c_0)A_R(u).
\]
The ideals are \(T\)-invariant, so the exact minimizing valuation is
\(T\)-invariant. Normalize it by \(A_R(u)=1\).

This \(u\) cannot be a cocharacter valuation in \(T\). For any centred
cocharacter competitor with bare discrepancy one, pure weights give
\(u(g)=m\) and \(u(I_R)=n\), hence
\(a=1+\lambda(c_0-1)\). Such competitors are lc places of the system;
\eqref{eq:A8} gives their volume at least
\(C^{-p}\vol(\tilde v)\). Their objective thus strictly exceeds
\(2^p\vol(\tilde v)\), while \(a(u_*)=1\) and the displayed volume
bound makes \(u_*\)'s objective at most that number. This contradicts
minimality if \(u\) were a cocharacter.

We have obtained a minimizer that is not a cocharacter of the old torus.
To contradict maximality, we must realize its joint graded ring on the
original \(W\) and preserve the system value \(n\), concentration
\eqref{eq:A7}, and the uniform comparison \eqref{eq:A8}.

Let \(R'=\gr_uR\). Stable degeneration for each sufficiently small
auxiliary klt parameter gives finite generation, normality and, after
removing boundaries, klt of \(R'\). Use the joint effective torus of
\(T\) and the new value torus. By
Lemma~\ref{a:canonical-specialization}, its canonical generator is
\((\Omega_{\tilde v})_u\); call its character \(\chi'\). It restricts
to \(\chi\) on \(T\), and \(u(\chi')=1\).
Choose the old core equations in \(R_{n\chi}\) for the exact-minimum
construction. Their next initials have old weight \(n\chi\) and
\(u\)-value \(n\), hence joint weight \(n\chi'\). The old and new
weights determine characters of the joint torus, so this is an equality
of characters, not only of their evaluations.

Let \(I_{R'}=(R'_{n\chi'})\). The core and normalized minimum remain
fixed as the auxiliary parameter tends to zero. Boundary specialization
and this pure weight therefore imply that \(I_{R'}^{1/n}\) is lc,
exactly as in the initial construction. The same auxiliary klt pairs give
a new uniform version of \eqref{eq:A8}, initially relative to grading
\(u\). Indeed at a centred lc place \(q\) of the new system, the core
term is at most zero; all perturbing boundary subtractions are
nonnegative, so the auxiliary discrepancy before common scaling is at
most
\[
                   t'(1+\lambda c_0)A_{R'}(q).
\]
At the minimizing grading \(u\) it equals \(t'a(u)\), which is at
least \(t'\). Formula~\eqref{eq:A3} cancels the parameter and yields
a sequence-uniform comparison with \(u\), since \(\lambda,c_0,p\) are
fixed.

On the old ring, \eqref{eq:A8} applies to the normalized lc place \(u\)
and gives \(u(f)\leq C\tilde v(f)\). On each new joint piece this says
that its new \(u\)-weight is at most \(C\) times its old
\(\tilde v\)-weight: use a homogeneous lift of that piece to \(R\).
We now flatten the two successive valuations on the actual original
\(W\), by Lemma~\ref{a:flattening}, and divide by \(1+\xi\).
This preserves \(D\), the iterated canonical generator, and discrepancy
one. The parameter \(\xi\) may be chosen individually for each member;
all bounds below use fixed factors.

Include a finite degree-\(n\) generating list for the actual ideal
\(I_n\) in the preserved tests. If an old value is greater than \(n\),
its normalized flattened value is still greater for small enough
\(\xi\). If the old value is \(n\), its first initial lies in
\(R_{n\chi}\), so the next value is at least \(n\). Include also a
degree-\(n\) lift attaining equality for this system minimum; such a
lift exists because \(R_{n\chi}\) is precisely that leading space.
The normalized flattened valuation therefore still has
\(\tilde v_{\mathrm{new}}(I_n)=n\).

The all-functions comparison in Lemma~\ref{a:flattening}, with
\(0<\xi<1\), bounds the new actual valuation below by
\(\tilde v/4\). Thus \eqref{eq:A7} persists. On each new joint weight,
its normalized value is
\[
 \frac{\text{old weight}+\xi\,\text{new }u\text{-weight}}{1+\xi}
       \geq\tfrac12\,\text{old weight}
       \geq\tfrac1{2C}\,\text{new }u\text{-weight}.
\]
Both grading valuations use the minimum rule, so the same comparison
holds on all elements of \(R'\). The new comparison relative to \(u\)
therefore gives \eqref{eq:A8} relative to the normalized flattened
grading, with a uniform enlarged constant. All required sequence
properties survive.

The joint torus has strictly larger rank. Every old occurring character
still occurs after taking the invariant associated graded, so \(T\)
continues to act faithfully. If the \(u\)-grading belonged to this old
torus, each nonzero element in an old weight space would have its
\(u\)-value determined by that character and one fixed linear
functional. Invariance and the minimum rule would then make \(u\)
itself a cocharacter valuation on \(R\), which was excluded above.
Thus the rank increases. Passing to a constant-rank subsequence gives
an admissible sequence of larger rank, contrary to its maximal choice.
This proves \eqref{eq:A12}.
\end{proof}

Fix this sequence from now on. The dual section
\[
 Q_T=\{\rho\in\Hom(M_T,\R):
          \rho\geq0\text{ on occurring weights},\ \rho(\chi)=1\}
\]
is compact, because \(\chi\) is interior and the occurring cone is
full-dimensional. It contains \(D\), and its point \(\tilde v\) obeys
\[
                    \tilde v+c(\tilde v-Q_T)\subseteq Q_T
\]
for a uniform \(c>0\). To see this, apply \eqref{eq:A8} first to
interior cocharacters \(\rho\). They are centred at the vertex and are
lc places of the pure system, of bare discrepancy one. Therefore
\(\rho(\mu)\leq C\tilde v(\mu)\) for every occurring character
\(\mu\). Continuity extends this to all \(Q_T\); take, for example,
any \(c>0\) with \(c(C-1)\leq1\).

Use the integral splitting
\[
                  M_T=\Z\chi\oplus M_0,\qquad M_0=\ker D.
\]
At each invariant prime divisor \(E\) on \(\Spec R\), the values
\(E(z^\beta)\) are integral linear functions of \(\beta\in M_0\).
These primes, with the data in \eqref{eq:A11}, test regularity of
eigenfunctions by normality: any pole divisor of an eigenfunction is
invariant, and a normal affine ring is the intersection of its
height-one valuation rings.

\subsection{Angular forms and section costs}
\label{a:angular}

Distinguish the quotient function \(f_h\) on \(S\) from the actual
rational function
\[
                     f_h^W=h/t^n\in k(X)=k(Z).
\]
Set
\[
 \theta=\frac{\omega}{(f_h^W)^{1/n}},\qquad
 V_m=h^{-m/n}\mathcal A_m\subset k(X)
                    \quad(m\in n\Z_{\geq0}).
\]
The form is fractional, interpreted by its \(n\)-th tensor power.
The section systems multiply and are nested, since multiplication by
\(h\) identifies \(V_m\subseteq V_{m+n}\). For
\(v=\tilde v|_{k(X)}\), the Gaussian cone formula gives
\(A_\theta(v)=0\). On \(0\ne F\in V_m\), write its cost as
\[
                              b(F)=m+v(F),
\]
which is the value of its actual degree-\(m\) section under
\(\tilde v\). The label \(m\) is retained when an element belongs
to more than one system.

More generally, for a quasi-monomial valuation \(w\) of \(k(X)\), the
Gauss extension assigning \(h\) value \(nB_*\) has
\[
 A_\Omega=A_\theta(w)+B_*,\qquad
             \text{section value}=w(F)+mB_*\quad(F\in V_m).
\]
Indeed its value on \(t\) is \(B_*-w(f_h^W)/n\); substitution in the
cone formula gives the displayed identities.

The full initial field for \(v\) is
\[
          k(X)_v=k(R)^D=k(S)(z^\beta:\beta\in M_0).
\]
This follows from Gaussianity of \(\tilde v\) in \(t\): its initial
field is the full initial field downstairs with one independent initial
variable of \(D\)-degree one. The leading \(V_m\)-spaces are the
leading degree-\(m\) pieces of \(R\) divided by \(\bar h^{m/n}\).
Apply \eqref{eq:A6} to the fractional identity
\[
                \Omega/h^{1/n}=\theta\wedge d\log t.
\]
For a basis of \(M_0\), the result is
\[
              \theta_v=\pm\frac{\omega_S\wedge
                                      d\log\mathbf z_0}{f_h^{1/n}}.
\]
In fact \(d\log\bar t=d\log z^\chi\) modulo differentials from
\(k(R)^D\), since both variables have \(D\)-degree one; wedging with
the full downstairs form kills their difference.

Finally let \(w\) be any quasi-monomial valuation centred on \(X\),
including one lying over \(z\). The general divisor
\(\frac1n(\Div_X f_h^W-nK_X)\) is lc, so \(A_\theta(w)\geq0\), and
\begin{equation}
 \begin{split}
 A_X(w)&=A_\theta(w)
             +\tfrac1n w(\Div_X f_h^W-nK_X),\\
 \tfrac1n w(\Div_X f_h^W-nK_X)
       &=\sup_{\substack{m\in n\Z_{>0}\\0\ne F\in V_m}}
                                      \frac{-w(F)}m.
 \end{split}
 \label{eq:A13}
\end{equation}
Effectivity of every degree-\(m\) section divisor gives the upper bound
for the supremum. Conversely, sufficiently divisible multiples of
\(-K_X\) are relatively free. Over the affine base choose a section
avoiding the centre of \(w\); its effective section divisor has value
zero there and attains the asserted value. This proves the equality
and completes the angular data used below.

Equation~\eqref{eq:A13} specifies the final task. A lower bound close
to zero for all positive-degree ratios \(w(F)/m\), together with small
\(A_\theta(w)\), gives small \(A_X(w)\). To use \(\epsilon\)-log
canonicity, we also need a uniformly bounded positive integer \(e\) for
which \(ew\) is a positive integral multiple of a prime-divisor valuation.
We will obtain this control from independent initial coordinates with
integral values and the section counts. The tower and section estimates
in the next four sections supply the information needed to construct
such a valuation on the original field.
